\begin{register}{H}{cpu\_gpio\_oen}{0x803}\label{regdesc:GPIO_OEN}
\regfield{(reserved)}{24}{8}{{0}}%
\regfield{CPU\_GPIO\_OEN[7]}{1}{7}{{0}}%
\regfield{CPU\_GPIO\_OEN[6]}{1}{6}{{0}}%
\regfield{CPU\_GPIO\_OEN[5]}{1}{5}{{0}}%
\regfield{CPU\_GPIO\_OEN[4]}{1}{4}{{0}}%
\regfield{CPU\_GPIO\_OEN[3]}{1}{3}{{0}}%
\regfield{CPU\_GPIO\_OEN[2]}{1}{2}{{0}}%
\regfield{CPU\_GPIO\_OEN[1]}{1}{1}{{0}}%
\regfield{CPU\_GPIO\_OEN[0]}{1}{0}{{0}}%
\reglabel{Reset}\regnewline%
\begin{regdesc}\begin{reglist}
\label{fielddesc:GPIO_OEN}\item [CPU\_GPIO\_OEN] 配置是否使能 GPIOn (n=0 \verb+~+ 7) 输出。CPU\_GPIO\_OEN[7:0] 分别对应章节 \textit{\nameref{mod:iomuxgpio}} 中表 \ref{tab:iomuxgpio-periph-signals-via-hp-gpio-matrix-cn} 里的 cpu\_gpio\_out\_oen[7:0] 输出使能信号。CPU\_GPIO\_OEN 的值与 cpu\_gpio\_out\_oen 的值对应。\\
此寄存器是 \hyperref[fielddesc:GPIO_OUT]{CPU\_GPIO\_OUT} 的使能寄存器。\\
0：关闭 GPIO 输出\\
1：使能 GPIO 输出\\
(R/W)
\end{reglist}\end{regdesc}
\end{register}

\begin{register}{H}{cpu\_gpio\_in}{0x804}\label{regdesc:GPIO_IN}
\regfield{(reserved)}{24}{8}{{0}}%
\regfield{CPU\_GPIO\_IN[7]}{1}{7}{{0}}%
\regfield{CPU\_GPIO\_IN[6]}{1}{6}{{0}}%
\regfield{CPU\_GPIO\_IN[5]}{1}{5}{{0}}%
\regfield{CPU\_GPIO\_IN[4]}{1}{4}{{0}}%
\regfield{CPU\_GPIO\_IN[3]}{1}{3}{{0}}%
\regfield{CPU\_GPIO\_IN[2]}{1}{2}{{0}}%
\regfield{CPU\_GPIO\_IN[1]}{1}{1}{{0}}%
\regfield{CPU\_GPIO\_IN[0]}{1}{0}{{0}}%
\reglabel{Reset}\regnewline%
\begin{regdesc}\begin{reglist}
\label{fielddesc:GPIO_IN}\item [CPU\_GPIO\_IN] 表示 SoC GPIOn (n=0 \verb+~+ 7) 的输入值（1 为高电平，0 为低电平）。\\
CPU\_GPIO\_IN[7:0] 分别对应章节 \textit{\nameref{mod:iomuxgpio}} 中表 \ref{tab:iomuxgpio-periph-signals-via-hp-gpio-matrix-cn} 里的 cpu\_gpio\_in[7:0] 输入信号。\\
CPU\_GPIO\_IN[7:0] 只能通过 GPIO 交换矩阵映射到 GPIO。详细描述请参考章节 \ref{sec:peripheral-input-via-gpio-matrix}。
\\
(RO)
\end{reglist}\end{regdesc}
\end{register}
\begin{register}{H}{cpu\_gpio\_out}{0x805}\label{regdesc:GPIO_OUT}
\regfield{(reserved)}{24}{8}{{0}}%
\regfield{CPU\_GPIO\_OUT[7]}{1}{7}{{0}}%
\regfield{CPU\_GPIO\_OUT[6]}{1}{6}{{0}}%
\regfield{CPU\_GPIO\_OUT[5]}{1}{5}{{0}}%
\regfield{CPU\_GPIO\_OUT[4]}{1}{4}{{0}}%
\regfield{CPU\_GPIO\_OUT[3]}{1}{3}{{0}}%
\regfield{CPU\_GPIO\_OUT[2]}{1}{2}{{0}}%
\regfield{CPU\_GPIO\_OUT[1]}{1}{1}{{0}}%
\regfield{CPU\_GPIO\_OUT[0]}{1}{0}{{0}}%
\reglabel{Reset}\regnewline%
\begin{regdesc}\begin{reglist}
\label{fielddesc:GPIO_OUT}\item [CPU\_GPIO\_OUT] 向 SoC GPIOn (n=0 \verb+~+ 7) 写输出值（1 为高电平，0 为低电平）。\hyperref[fielddesc:GPIO_OEN]{CPU\_GPIO\_OEN} 置位时，写输出值才有效。\\
CPU\_GPIO\_OUT[7:0] 分别对应章节 \textit{\nameref{mod:iomuxgpio}} 中表 \ref{tab:iomuxgpio-periph-signals-via-hp-gpio-matrix-cn} 里的 cpu\_gpio\_out[7:0] 输出信号。\\
CPU\_GPIO\_OUT[7:0] 只能通过 GPIO 交换矩阵映射到 GPIO。详细描述请参考章节 \ref{sec:peripheral-output-via-gpio-matrix}。
\\
(R/W)
\end{reglist}\end{regdesc}
\end{register}
